% =====================================================================
%  Urban Innovation Challenge — Research Report
%  Apple Developer Academy @ BINUS, Cohort 9
%
%  Compiled on OVERLEAF (pdfLaTeX). No local build required.
%  Prose language: Bahasa Indonesia. Technical terms may be English.
%
%  NOTE: babel option is 'indonesian' (NOT 'bahasa' — 'bahasa' errors
%  on modern babel).
% =====================================================================
\documentclass[11pt,a4paper]{report}

\usepackage[utf8]{inputenc}
\usepackage[T1]{fontenc}
\usepackage[indonesian]{babel}   % <-- 'indonesian', never 'bahasa'
\usepackage{graphicx}
\usepackage{hyperref}
\usepackage{csquotes}

% --- Bibliography (uncomment when refs.bib has entries) ---------------
% \usepackage[backend=biber,style=apa]{biblatex}
% \addbibresource{refs.bib}

% =====================================================================
\title{%
  Pengalaman Berjalan Kaki di Koridor Sudirman\\[0.4em]
  \large Urban Innovation Challenge --- Challenge Based Learning}
\author{%
  Tim \texttt{[NAMA TIM EKSTERNAL --- TBD]}\\
  Apple Developer Academy @ BINUS, Cohort 9}
\date{\today}

\begin{document}
\maketitle
\tableofcontents

% =====================================================================
\chapter{Pendahuluan}            % BAB I
% TODO: Latar belakang, Big Idea ("Urban Living Experience"),
% pernyataan tantangan (sementara), dan ruang lingkup (koridor Sudirman
% + jalan-jalan di belakangnya, berbasis nama jalan).
Bab ini akan menjelaskan latar belakang dan ruang lingkup penelitian. \textit{(Placeholder.)}

% =====================================================================
\chapter{Tinjauan Pustaka}       % BAB II
% TODO: Literature review dari desk-research/. Gunakan sitasi dari refs.bib.
Bab ini akan merangkum studi terkait keterjalankakian (\textit{walkability}). \textit{(Placeholder.)}

% =====================================================================
\chapter{Metodologi}             % BAB III
% TODO: Metode lapangan (audit trotoar, observasi/perhitungan, intercept,
% wawancara ahli), instrumen, dan etika penelitian.
Bab ini akan menjelaskan metode pengumpulan dan analisis data. \textit{(Placeholder.)}

% =====================================================================
\chapter{Hasil \& Kesimpulan}    % BAB IV
% TODO: ISI HANYA SETELAH DATA LAPANGAN TERKUMPUL. Jangan menulis temuan
% atau kesimpulan sebelum data nyata tersedia.
Bab ini akan memuat hasil dan kesimpulan setelah kerja lapangan. \textit{(Placeholder.)}

% --- Bibliography (uncomment together with biblatex above) -----------
% \printbibliography

\end{document}
